\chapter{State of the Art}
\label{chap:state-of-the-art}

Improving the efficiency of AI inference is not a single-layer problem. A system can reduce the amount of computation or memory required by a model, make better use of an already allocated accelerator, control the operating point of the hardware, or improve how devices are shared across a cluster. These approaches are complementary, but they optimize different quantities and introduce different constraints. A higher throughput inside one inference server, for example, does not necessarily eliminate unused capacity elsewhere in the cluster. Similarly, a smaller model does not guarantee that the platform will allocate a smaller resource to it.

This chapter reviews representative techniques at four levels: model, inference runtime, hardware control, and cluster orchestration. The purpose is not to provide an exhaustive survey, but to locate the infrastructure problem addressed by this thesis within the broader work on efficient AI inference.

\section{Model-Level Optimizations}
\label{sec:sota-model-level}

Quantization reduces the numerical precision used to represent model weights and, in some cases, activations. Replacing floating-point values with lower-precision representations can reduce the model footprint and memory traffic, and can accelerate matrix operations when suitable hardware kernels are available. SmoothQuant, for example, enables post-training W8A8 quantization of large language models by moving part of the quantization difficulty from activation outliers to the weights through an offline, mathematically equivalent transformation. The authors report approximately half the memory use and speedups of up to 1.56$\times$ for the evaluated models and hardware~\cite{xiao2023smoothquant}.

The advantage is obtained by changing the model representation rather than the allocation policy. Its effectiveness therefore depends on the selected precision, the sensitivity of the model, calibration data, and the availability of efficient kernels on the target accelerator. More aggressive quantization can degrade output quality, while conversions and operations that remain at higher precision can reduce the expected speedup. Furthermore, even a quantized model may still be assigned an entire GPU by the surrounding platform. Quantization reduces the demand created by one inference execution, but does not by itself decide how the released capacity should be shared with other workloads.

\section{Inference-Runtime Optimizations}
\label{sec:sota-runtime-level}

\subsection{Batching and Request Scheduling}
\label{subsec:sota-batching}

Batching combines several inputs into one model execution. This can increase arithmetic intensity and amortize kernel-launch and data-transfer overheads, thereby raising throughput on an allocated GPU. NVIDIA Triton Inference Server implements dynamic batching by collecting compatible requests and forming batches at runtime. Its scheduler can wait for a configurable interval to obtain a preferred batch size, making the throughput--latency trade-off explicit~\cite{nvidia2026tritonbatching}.

The principal limitation is that requests may spend additional time in the queue while a batch is formed. Large batches can also increase execution time, memory consumption, and head-of-line blocking, especially when inputs have heterogeneous shapes or sequence lengths. Under light or irregular traffic, there may not be enough simultaneous requests to form an efficient batch. Consequently, batching must be tuned against latency objectives rather than enabled under the assumption that the largest possible batch is always preferable.

Autoregressive language models make this problem more complex because requests generate different numbers of tokens and therefore finish at different times. Orca addresses this limitation with iteration-level scheduling: instead of keeping a request in a fixed batch until completion, the runtime can reconsider the batch at each decoding iteration and admit new work when capacity becomes available~\cite{yu2022orca}. This improves utilization for persistent generative-model services, but requires control of the execution loop and a specialized serving runtime. It is less directly applicable when each inference task is packaged as an independent, short-lived process rather than submitted as a request to a shared server.

\subsection{Memory Management}
\label{subsec:sota-memory-management}

For large language models, GPU memory is consumed not only by model weights but also by the key--value (KV) cache maintained for active sequences. The cache has a variable lifetime and size, so contiguous reservation can cause internal fragmentation and limit the number of requests that can execute concurrently. vLLM introduces PagedAttention, which stores KV-cache data in fixed-size, non-contiguous blocks and allows blocks to be shared when requests have common content. The original evaluation reports a 2--4$\times$ throughput improvement over the considered serving baselines at a comparable latency level~\cite{kwon2023pagedattention}.

This optimization is particularly valuable for high-concurrency, autoregressive serving, but it is not a general solution to cluster-level GPU underutilization. Its benefits depend on the presence of KV-cache pressure, the request-length distribution, and sufficient concurrency. It also couples deployment to a runtime that supports the required models and decoding features. More efficient memory management allows an inference server to host more active sequences; it does not determine whether several independent applications should receive the same physical accelerator.

Systems such as Clockwork extend runtime control to model loading, placement, and request scheduling. Clockwork centralizes these decisions to provide predictable latency while sharing GPUs among many deployed DNN models~\cite{gujarati2020clockwork}. The resulting control can be effective for an integrated serving platform, but requires applications to use that platform and transfers substantial responsibility to a specialized controller.

\section{Hardware-Control Optimizations}
\label{sec:sota-hardware-level}

Efficiency can also be improved by changing the hardware operating point. Dynamic voltage and frequency scaling and power capping reduce instantaneous power when the workload has performance margin. The $\mu$-Serve system jointly optimizes model multiplexing and GPU frequency for multi-model serving. On the production workloads evaluated by its authors, it reduced power by up to 61\% without violating the considered service-level objectives~\cite{qiu2024muserve}.

Power control exposes another trade-off: a lower power level can lengthen execution, so lower watts do not necessarily imply lower energy per completed request. A controller needs reliable workload profiles and must react when arrival rates or latency requirements change. Conservative settings leave energy-saving opportunities unused, whereas aggressive settings risk increasing tail latency. Moreover, frequency control does not solve allocation fragmentation: a GPU can operate efficiently while still being reserved by a workload that uses only a small part of its capacity.

Hardware concurrency and partitioning mechanisms provide a different opportunity by allowing multiple workloads to make progress on one device. They can improve consolidation, but temporal sharing may introduce interference and spatial isolation may reduce flexibility or leave unusable fragments. The appropriate mechanism therefore depends on memory capacity, performance predictability, security, and fault-containment requirements. These mechanisms are described in detail in Chapter~\ref{chap:background}; at the state-of-the-art level, the important observation is that hardware sharing exposes capacity, while a higher-level policy must still decide which workloads may safely use it.

\section{Cluster-Orchestration Optimizations}
\label{sec:sota-platform-level}

At cluster level, the objective changes from accelerating one model invocation to placing and isolating multiple workloads across the available devices. HAMi (Heterogeneous AI Computing Virtualization Middleware) is an open-source Kubernetes project that addresses this layer. For NVIDIA GPUs, it allows a Pod to request a quantity of device memory and a percentage of compute capacity instead of treating every request as an exclusive whole-device allocation. Its scheduler extender maintains a device-level view of remaining resources and can apply bin-packing or spreading policies. A device plugin exposes the resources to Kubernetes, while an injected library intercepts CUDA operations to enforce memory limits and throttle compute usage inside the container~\cite{hami2026virtualization}.

This design can improve density when several workloads have modest and compatible requirements. In particular, bin-packing can concentrate fractional allocations on fewer devices and leave other GPUs available for larger requests. HAMi also provides a common management approach for heterogeneous accelerators, although the supported sharing granularity and isolation capabilities vary by device vendor~\cite{hami2026faq}.

HAMi illustrates both the value and the limits of a platform-level solution. Resource requests still need to represent application demand accurately: an undersized quota can reduce performance, while an oversized one recreates waste. Consolidating workloads can introduce contention and variable latency, and CUDA API interception provides software-enforced limits rather than hardware fault isolation. The additional webhook, scheduler, device-plugin, and runtime components also increase operational complexity and compatibility requirements. Finally, HAMi does not automatically optimize model precision, batch size, or KV-cache management; those decisions remain at the model and serving-runtime levels.

\section{Research Gap and Thesis Positioning}
\label{sec:sota-positioning}

Much of the influential inference-systems literature reviewed above is organized around long-running services. Triton, Orca, vLLM, and Clockwork assume that models remain available while a stream of requests is batched, scheduled, or multiplexed. This is a natural design for popular LLM endpoints, because model loading is expensive and sustained concurrency creates opportunities for continuous batching and cache reuse. It is less representative of small, finite inference tasks that arrive as separate jobs, execute a bounded amount of work, and then release their resources. For such workloads, Pod admission, device selection, startup overhead, and cluster fragmentation can matter as much as request scheduling inside the model server.

The work in this thesis therefore operates mainly at the boundary between hardware resource management and Kubernetes orchestration, rather than modifying model architecture or implementing a new inference runtime. It investigates whether hardware-isolated GPU subdivisions can be exposed through Kubernetes Dynamic Resource Allocation (DRA), so that device selection can be reconsidered for successive finite workloads instead of being fixed permanently for a long-running service. The current batch mechanism is first validated with finite GPU jobs; extending it to heterogeneous inference jobs remains a distinct experimental step. The research question is whether integrating DRA with admission and allocation policy can improve the use of already available GPU capacity relative to credible static baselines. Model-level and runtime techniques such as quantization, batching, and PagedAttention remain complementary and could be applied within an allocated resource, but they do not replace this cluster-level decision.
